\section{The modular arm}
\label{app:modular}

The modular arm infers the reduced model one block at a time on the dataset of
\cref{sec:reduced}, and sets the result beside the joint posterior. The result is a cut
distribution, a different target from the joint posterior
\cite{plummer2015cuts,jacob2017better}. We fixed its design, its closeness rule and its
scoring in a pre-registration tagged before any modular draw existed (\texttt{prereg/kim-116}
in the repository).

\paragraph{Two variants.} Both have two stages.
\begin{itemize}
  \item \emph{The factor cut.} Stage one is the transcript factor alone, over all
  \pn{kim92.ncoord} coordinates, with ENO, FBA and $\sigb$ under their priors. It still
  solves the mean-field ODE, because the transcript factor reads its rate constants from
  that solve.
  \item \emph{The data plug-in.} Stage one takes its rate constants from the observed ATP
  and GTP panel, with CTP and UTP at their chemostat values, and needs no ODE solve.
\end{itemize}
Stage one uses the joint's chain design in both variants.
In stage two, both variants carry \pn{kim116.K} stage-one draws forward. For each, a
No-U-Turn chain fits ENO, FBA and $\sigb$ to the bulk factor, with the promoters and the decay
rate fixed at that draw. It adapts for \pn{kim116.inner} iterations, then keeps the next
\pn{kim116.innerkeep} draws.

\paragraph{Gates.} Stage one had to reach $\Rhat \le \pn{kim92.rhatgate}$ and a bulk
effective sample size of at least \pn{kim92.essgate} on the promoters and the decay rate,
with one extension
allowed. Both variants needed the extension, and both then passed. In the posterior
predictive check of \cref{sec:reduced}, \pn{kim116.ppcout} of the \pn{kim92.ppcrows} bulk
rows fall outside the central 99\% band for each variant, against the joint's
\pn{kim92.ppcout}.

\paragraph{Cost.} The factor cut's stage one costs the same as the joint's chains, because
it solves the same ODE over the same coordinates. Stage two comes on top, at
\pn{kim116.stagetwo}\% of the cut's total. So the cut costs \pn{kim116.cutjob} times the
joint's own posterior jobs. The plug-in's stage one costs \pn{kim116.plugstageone} times the
joint's chains, and stage two is \pn{kim116.stagetwoplug}\% of its total, so the plug-in costs
\pn{kim116.plugjob} times the joint's jobs. Against the whole joint run of
\cref{tab:reducedcost}, the pre-registered comparator, the ratios are \pn{kim116.cutt1} and
\pn{kim116.plugt1}. Every arm also shares the joint's dataset, preparation, probe and profile; the last three
cost \pn{kim116.shared} CPU-hours, paid once. The size of stage two was fixed in the
pre-registration and never tuned for cost, so the plug-in's ratio belongs to this design and
is not a floor.

\paragraph{Closeness to the joint.} The factor cut is close by the pre-registered rule. Its
largest shift in a mean is \pn{kim116.cutshift} joint SD, and its largest change in
contraction is \pn{kim116.cutcontr}. The plug-in fails the rule: the GAPD promoter shifts by
\pn{kim116.plugshift} joint SD, with a Monte Carlo SE of \pn{kim116.plugshiftse}. The bias is
systematic: \pn{kim116.plugup} of the \pn{kim92.npromoter} promoters shift up, and the mean
shift over all \pn{kim92.npromoter} is \pn{kim116.plugmean} SD. Its largest change in contraction is
\pn{kim116.plugcontr}.

\paragraph{Distance from the truth.} The joint has no exact reference on the reduced model,
so each variant's distance from the truth sits beside its distance from the joint. Both
variants hold the truth inside \pn{kim116.truthin} of the \pn{kim92.ncoord} central 90\%
intervals, where the joint holds \pn{kim92.truthin}. The cut's promoters have a mean
standardised error of \pn{kim116.cutpromz}, the plug-in's \pn{kim116.plugpromz}, and the
joint's \pn{kim92.promz}. The plug-in's value is not a sign that it sits nearer the truth.
We read it as the plug-in's own upward bias laid over the same downward offset, and we have
not tested that reading.

\paragraph{Why the plug-in shifts (not tested).} At the truth, the plug-in's rate constants
differ from the mean-field route's by a median of \pn{kim116.dlogkmed} on the log scale, and
by at most \pn{kim116.dlogkmax}. That is an unsigned size; it does not by itself show a bias
in one direction. The profile names four candidate sources. The panel's noise is treated as the true
state, an errors-in-variables cut. The offset $\delta$ is applied. A save stands in for the
switch one second before it. And the volume and the pools outside the panel come from the
joint's reference solve. Which of them moves the promoters up was not tested.

\paragraph{Scored predictions.} Only the plug-in was scored, as the pre-registration fixed,
on the targets, sizes, truth and prior arm of \cref{sec:reducedpredict}. On the knockdown's
effect $\Delta Q$ it scores as the joint does, with a CRPS of \pn{kim116.dqcrpsplug} against
\pn{kim116.dqcrpsjoint}. On the pools and transcripts at the later horizon it is within a few
points of the joint. On $Q$ and $Q_{\mathrm{kd}}$ its median sits about
\pn{kim116.plugqhigh}\% above the truth. Its CRPS on $Q$ is \pn{kim116.plugqcrps}, about
twice the joint's \pn{kim116.jointqcrps}, though its skill against the prior stays at
\pn{kim116.plugskill}.
